\section{Validating vertex links by an incidence census}
\label{sec:vertex-link-census}

The shelling measurements expose a cost shared by discovery and independent
verification: repeatedly reconstructing and validating finite manifold
geometry. The previous validator built four union--find structures, for
vertices, signed edges, directed edges and vertex-link edges. The latter two
can be removed without dropping a validation condition. Signed edge classes
already determine all vertex-link vertices, while triangle-side incidence
determines the number of vertex-link edges. The new census retains the
existing source checks, rejection rules, matching equations and normal-surface
proof formats. It also applies to supplied triangulations whose source is not
a knot diagram; no provenance assumption is used to bypass manifold checks.

\subsection{Oriented ends and the two ends of a loop}

A local tetrahedron corner contributes a triangle to the link of its vertex.
A vertex of that triangle corresponds to a directed edge leaving the corner.
Under a face gluing, these directed edge occurrences are identified exactly
when they represent the same end of the same global edge. The signed edge
union already records whether each local ascending edge agrees with its
chosen representative. Its consistency check rejects an edge identified with
its own reverse. Therefore every global edge has precisely two distinct
oriented-end classes, each attached to its corresponding global vertex.

This remains true for a loop edge whose two endpoints have the same global
vertex. The ends are different vertices of that vertex's link, even though
they contribute to the same link. Counting a loop once would undercount its
link vertices and invalidate the Euler test. The implementation increments
the endpoint census twice for every global edge, including when the endpoint
labels coincide. Reversing a representative's orientation only exchanges its
two endpoints, so their incidence counts are independent of that choice.
The existing endpoint map, which is also returned for later normal geometry,
is reused for this computation.

Equivalently, a local directed edge can be labelled by its signed global edge
representative and one orientation bit. Reconstructing all such labels with
a second union--find computes an equivalence relation already contained in
the signed edge structure. The implementation needs only their counts at
each vertex, so it does not even materialize the doubled collection of ends.

\subsection{Counting link edges without reconstructing them}

For a global vertex $v$, let $F_v$ be its number of tetrahedral corners, let
$V_v$ be the number of incident global edge ends, and let $B_v$ count unpaired
triangle-side occurrences in its link. Every link triangle has three sides.
An internal link edge pairs exactly two side occurrences, and a boundary
link edge has exactly one. Hence, if $I_v$ and $E_v$ denote its numbers of
internal and total edges,
\[
 3F_v=2I_v+B_v,\qquad E_v=I_v+B_v=\frac{3F_v+B_v}{2}.
\]
It follows that
\[
 2\chi(\operatorname{Lk}v)=2V_v-F_v-B_v.
\]
Reciprocity and the prohibition on pairing a face with itself are checked
before this identity is used. A link side is paired across its one containing
tetrahedral face, so there is no chain that merges three or more side
occurrences into one link edge. Pairing different sides of one link triangle
is allowed and still contributes two occurrences, as the identity requires.
Boundary counts use local incidences with multiplicity: a boundary triangle
whose vertices become identified can contribute more than once to one $B_v$.

All earlier manifold prerequisites remain active. Face-connectedness and
orientability are checked. The signed edge union rejects reversal, and each
global edge link must have zero or two unpaired endpoints; its incidence
graph is connected by construction and has degree at most two. These checks
make vertex links connected surfaces. The finite-vertex test then requires
$2V_v-F_v-B_v=4$ when $B_v=0$, and $2V_v-F_v-B_v=2$ otherwise. These are the
same sphere and disc Euler conditions as before, evaluated by a different
cell census. The final boundary checks still require exactly one torus.
An ideal torus vertex has link Euler characteristic zero and is still rejected.
Thus the optimization does not accept ideal or singular vertices by assuming
that an authenticated source has already made them harmless.

\subsection{Operation counts, storage and preserved outputs}

Three arrays indexed by the original vertex roots accumulate $F_v$, $B_v$
and $V_v$. They have $4N$ entries each for $N$ tetrahedra. The first counts
local corners, the second counts corners of unpaired faces, and the third
counts the two endpoints of every global edge. Once the existing global
classes and endpoint map have been obtained, this census requires $O(N)$
indexed integer operations and $O(N)$ words. Every counter has
$O(\log(N+1))$ bits. No normal sheet or integer height is inspected or expanded.
The earlier approach performed additional union--find work on both directed
edges and triangle sides, followed by per-link sets of representatives.

Let $P$ be the number of paired tetrahedral face pairs, and let $b$ be the
number of unpaired faces. Before the unchanged boundary-component union,
the old validator allocates $34N$ union--find nodes, while the new validator
allocates $10N$. Including the boundary-component structure gives $34N+b$
and $10N+b$ nodes. Every paired face formerly induced 15 unions: three on
vertices, three on link sides, three on signed edges and six on directed
edges. It now induces six. If $e_\partial$ is the number of global boundary
edges, the complete union counts are therefore $15P+e_\partial$ and
$6P+e_\partial$. The measured family below checks these exact counts:
\begin{center}
\small
\begin{tabular}{rrrrr}
\toprule
Tetrahedra & Old nodes & New nodes & Old unions & New unions \\
\midrule
20 & 688 & 208 & 552 & 228 \\
160 & 5,504 & 1,664 & 4,416 & 1,824 \\
640 & 22,016 & 6,656 & 17,664 & 7,296 \\
2,560 & 88,064 & 26,624 & 70,656 & 29,184 \\
\bottomrule
\end{tabular}

\end{center}
Each union--find node has parent, size and parity entries. Removing the two
structures eliminates $72N$ list entries, while the three new census arrays
add $12N$. This describes those named structures, not the peak memory of the
entire Python process. The old per-link dictionaries and sets also disappear;
the endpoint map was already required by the returned geometry. Caching each
signed edge's root and parity together avoids finding it again at export.

The complete validator still has the same near-linear indexed union--find
bound and linear-size combinatorial output. The usual dictionary and set
collision qualifications still apply to its remaining maps. The improvement
removes redundant equivalence computations and makes the additional link
census linear; it does not change the worst-case complexity class of general
unknot recognition. Vertex labels, signed edge labels, matching equations,
boundary incidence, endpoint representatives and all other returned fields
are preserved exactly. Existing guard callbacks stay at the same locations,
so work budgets and certificate-selection behavior remain unchanged even
when actual execution becomes faster.

\subsection{Exact compatibility and independent manifold controls}

The audit loads the geometry validator from revision \code{635465da5} and
records its complete source hash. It compares every returned field, rejection
message and callback count, not just the final recognition verdict. On each
of the 84 source diagrams it compares the current pulling exterior, its
shelling reduction and the legacy centred exterior. Nine layered solid tori,
from one to 256 tetrahedra, bring the exact geometry comparison to 261 inputs.
These include one-vertex triangulations with loop edges and canonical
exteriors with interior vertices.

All 252 published native-search records are reproduced exactly: 84 sources
each with shellings off, shellings on, and shellings plus planar capping.
The 47 distinct positive certificates are unchanged and replay successfully.
This includes the previous work-cap outcomes and all stage summaries. Lower
CPU cost does not change the deterministic work allowance, whose guard-count
schedule is preserved. The 19 ordinary pipeline controls still finish through
earlier recognizers, with no native attempt, and their verdicts and methods
are unchanged.

The audit also generates 12,000 reciprocal face pairings on one to eight
tetrahedra, alternating unrestricted gluing permutations with orientation-
compatible ones. The old and new validators agree on all prepared data or
rejection messages and on every callback count. Regina independently checks
the supported-manifold classification for all 12,000: connected, orientable,
valid, finite, and with exactly one torus boundary. Both native implementations
accept 684 cases. Among the rejections, 1,985 fail the vertex-link condition,
2,897 are nonorientable, 4,877 are not face-connected, 378 have no finite
boundary and 1,179 have a nontorus boundary. Seed, generator code, classification
counts and a hash of the ordered generated records are retained for reproduction.

New focused controls explicitly reconstruct local link triangles, directed
ends and paired sides by graph traversal. Thirty relabelled fixtures compare
that literal census with accepted geometry, including both canonical
subdivisions and one-vertex layered solid tori. A fixed two-tetrahedron ideal
figure-eight exterior has a torus vertex link and must still fail the finite
vertex test. Cancellation and input immutability are also checked. All 42
focused tests pass in 3.655 seconds, and all 1,232 maintained tests pass in
216.855 seconds. The final source and random-pairing audit takes 55.020 seconds
and partly overlaps the full suite; these elapsed values are validation costs,
not performance comparisons. The audit and timing records pin 510 source files.

\subsection{Complete recognition and proof-replay measurements}

The isolated comparison replaces only the geometry validator, loading the
old implementation before timing. Both arms use the same selection wrapper.
Each timed call runs complete recognition with earlier shortcuts disabled,
the native stage enabled, a twelve-second allowance and a 20,000-object
fallback allowance, followed by an independent replay and hash of any native
proof. Native misses continue through the reduced-Khovanov fallback. Cases
marked ``plain'' disable shellings; cases marked ``shell'' enable them. The
old/new comparison within each row keeps that choice fixed.

Five randomized rounds of old, old A/A, new and new A/A calls cover eleven
pipelines after one warm-up round each: 220 measured calls and 44 warm-ups
in 15.901 seconds. A new-arm A/A ratio of 1.053 on the trefoil motivates a
nine-round repeat. It completes 396 measured calls and 44 warm-ups in
29.043 seconds. All 616 measured calls and 88 warm-ups complete correctly.
Within each case, verdict, method, native work count and certificate hash
are identical across all four arms and every round. No tests, audits, builds
or commits overlap either timing run.
\begin{center}
\small
\begin{tabular}{lrrrrr}
\toprule
Pipeline & Old ms & New ms & Old/new & Old A/A & New A/A \\
\midrule
Circle 9, plain & 123.147 & 93.930 & 1.313 & 1.007 & 1.000 \\
Circle 9, shell & 109.032 & 75.259 & 1.425 & 1.007 & 1.004 \\
Random 29, plain & 31.734 & 23.707 & 1.338 & 0.990 & 0.997 \\
Random 29, shell & 57.294 & 41.883 & 1.389 & 1.020 & 0.986 \\
Original disc, plain & 72.279 & 61.827 & 1.158 & 1.002 & 0.999 \\
Original disc, shell & 69.935 & 50.392 & 1.390 & 1.009 & 1.002 \\
Random 142, shell & 141.855 & 114.957 & 1.238 & 0.993 & 0.995 \\
Random 1, shell & 44.022 & 32.310 & 1.357 & 1.002 & 0.988 \\
Random 81, shell & 68.527 & 49.577 & 1.385 & 1.001 & 1.008 \\
trefoil, shell & 24.406 & 20.977 & 1.164 & 0.999 & 0.995 \\
Three-crossing, shell & 39.935 & 28.772 & 1.393 & 1.004 & 1.007 \\
\bottomrule
\end{tabular}

\end{center}
All eleven final paired old/new medians favor the new census, from 1.158 to
1.425. Final A/A medians range from 0.986 to 1.020. Plain raw witnesses improve
by factors of 1.313 and 1.338; the plain optimized-disc fixture improves by
1.158. Shelling pipelines improve by 1.238--1.425 on the positive fixtures,
and the trefoil's complete negative pipeline improves by 1.164. The benefit
comes from repeated geometry construction and verification, while unchanged
source authentication, cocycle extraction and optimization limit the total
gain. Guard counts stay unchanged because guards do not count every union or
array operation.

Absolute timing can drift between cases. In particular, the plain circle
row was measured during a slower interval than the shelling-circle row.
These rows must not be used as an off/on comparison. A separate randomized
experiment measures that choice under the new census, using the preceding
shelling driver unchanged. Its 200 measured calls and 40 warm-ups all finish
correctly in 10.602 seconds, with A/A medians from 0.989 to 1.048. Nine of ten
pipelines still favor shellings off, with paired off/on ratios from 0.577 to
0.857. The original optimized-disc fixture now favors shellings by a factor
of 1.245, but this one favorable case does not justify a general default.
All four added native corpus positives still take longer with shellings than
with their earlier Khovanov fallback in this full-verification experiment.
The shelling option therefore remains disabled by default.

The incidence census is enabled for all finite-manifold geometry preparation;
there is no new recognition option or certificate schema. Its public
validation contract is unchanged, and the ordinary pipeline corpus reaches
no native stage, so these measurements do not establish a default-portfolio
speedup on that corpus. They establish a reduction in geometry work and
measured gains in complete native pipelines, including the independent proof
cost that defeated the previous simplification experiment.

Implementation revision \code{16e918cc7} follows a fresh merge check against
main. Reproduce the validation with the driver
\code{normal\_orbit\_research.link\_census} in \code{audit} mode, and the
paired timing comparison in \code{recognize --rounds 9} mode. The separate
shelling-choice run uses \code{normal\_orbit\_research.shellings} in
\code{recognize --rounds 5} mode. All measured source proofs and discovery
families are unchanged. In particular, this optimization does not supply the
missing general subexponential or quasipolynomial discovery bound. The incoming
unknot inventory and \code{openai/math} revision \code{fd4aeeb2e} remain unchanged.


